\documentclass[11pt]{article}
\usepackage[a4paper,margin=1in]{geometry}
\usepackage{fontspec}
\usepackage[boldfont=false]{xeCJK}
\setCJKmainfont{FandolSong-Regular.otf}
\usepackage{amsmath}
\usepackage{amssymb}
\usepackage{array}
\usepackage{booktabs}
\usepackage{graphicx}
\usepackage{adjustbox}
\usepackage{enumitem}
\setlist[itemize]{topsep=2pt,partopsep=0pt,parsep=0pt,itemsep=2pt,leftmargin=1.4em}
\setlist[enumerate]{topsep=2pt,partopsep=0pt,parsep=0pt,itemsep=2pt,leftmargin=1.6em}
\usepackage{parskip}
\usepackage{fvextra}
\fvset{breaklines=true,breakanywhere=true,fontsize=\small}
\setlength{\parindent}{0pt}

\begin{document}

\begin{center}
  {\LARGE\bfseries Finance and accounting}\\[0.35em]
  {\large A Level Business}
\end{center}
\vspace{0.6em}

\section*{Why a business needs finance}

\textbf{Business finance} 企业资金 is the money a firm needs to start, run and grow. A new firm needs money to buy \textbf{premises} 经营场所 and equipment; a growing firm needs money for new projects; and every firm needs money to pay day-to-day bills.

\section*{Capital and revenue expenditure}

Spending falls into two types:

\begin{itemize}
\item \textbf{capital expenditure} 资本性支出 — money spent on \textbf{fixed assets} 固定资产 that last a long time, such as machines, vehicles and buildings.

\item \textbf{revenue expenditure} 收益性支出 — money spent on day-to-day running, such as wages, rent and raw materials.

\end{itemize}

\section*{Working capital}

\textbf{Working capital} 营运资本 is the money a firm has for its daily needs. It is found from the \textbf{balance sheet} 资产负债表:

\[
\text{working capital} = \text{current assets} - \text{current liabilities}
\]

Here \textbf{current assets} 流动资产 are things that become cash within a year (cash, inventory, money owed by customers), and \textbf{current liabilities} 流动负债 are debts due within a year. Too little working capital is dangerous — the firm may not be able to pay its bills.

\section*{Cash is not the same as profit}

This is a key idea.

\begin{itemize}
\item \textbf{profit} 利润 is what is left when total costs are taken from total revenue over a period.

\item \textbf{cash} 现金 is the money the firm actually has right now to spend.

\end{itemize}

A firm can be \textbf{profitable} but still run out of cash — for example, if it sells goods but customers have not paid yet. Many firms fail not from low profit, but from running out of cash.

\section*{Sources of finance}

Finance can come from \textbf{inside} the business (\textbf{internal sources} 内部来源) or from \textbf{outside} (\textbf{external sources} 外部来源). It can also be \textbf{short-term} 短期 (paid back within a year) or \textbf{long-term} 长期.

\begin{center}
\adjustbox{max width=\linewidth}{%
\begin{tabular}{lll}
\toprule
\textbf{Source} & \textbf{Type} & \textbf{What it is} \\
\midrule
\textbf{retained profit} 留存利润 & internal, long & profit kept in the business instead of paid out \\
\textbf{sale of assets} 出售资产 & internal & selling items the firm no longer needs \\
\textbf{share capital} 股本 & external, long & money raised by selling shares \\
\textbf{loan} 贷款 & external, long & borrowed money repaid with interest \\
\textbf{debenture} 公司债券 & external, long & a long-term loan certificate sold by a company \\
\textbf{overdraft} 透支 & external, short & the bank lets the account go below zero for a short time \\
\textbf{leasing} 租赁 & external & renting an asset instead of buying it \\
\textbf{hire purchase} 分期付款 & external & buying an asset by paying in instalments \\
\textbf{trade credit} 商业信用 & external, short & paying a supplier later (e.g. after 30 days) \\
\textbf{venture capital} 风险资本 & external, long & money from investors into risky young firms \\
\textbf{crowdfunding} 众筹 & external & small amounts raised from many people online \\
\textbf{microfinance} 小额信贷 & external & small loans for people with little access to banks \\
\bottomrule
\end{tabular}}
\end{center}

\section*{Choosing a source of finance}

There is no single best source. The choice depends on:

\begin{itemize}
\item the \textbf{amount} needed — large amounts may need shares or a long loan.

\item the \textbf{purpose} — long-term assets should use long-term finance.

\item the \textbf{cost} — interest and fees differ between sources.

\item the \textbf{legal structure} — only a company can sell shares.

\item how much the firm already owes, and whether it can offer \textbf{collateral} 抵押品 (an asset the lender can take if the loan is not repaid).

\end{itemize}

\section*{Cash-flow forecasts}

A \textbf{cash-flow forecast} 现金流量预测 predicts the cash coming in and going out each month. It helps a firm spot a shortage early.

\begin{itemize}
\item \textbf{cash inflows} 现金流入 — money coming in, mostly from sales.

\item \textbf{cash outflows} 现金流出 — money going out, such as wages and rent.

\item \textbf{net cash flow} 净现金流 — inflows minus outflows for the period.

\end{itemize}

\[
\text{net cash flow} = \text{cash inflows} - \text{cash outflows}
\]

The \textbf{opening balance} 期初余额 is the cash at the start of the month. Adding net cash flow gives the \textbf{closing balance} 期末余额, which becomes next month's opening balance.

\section*{Cash-flow problems and how to fix them}

Common causes of cash-flow problems are: too many sales on credit, holding too much stock, buying too many assets at once, and \textbf{overtrading} 过度交易 (growing too fast for the cash available).

Ways to improve cash flow:

\begin{itemize}
\item bring cash in sooner — ask customers to pay faster, or offer discounts for quick payment.

\item delay cash going out — agree longer trade credit with suppliers.

\item arrange an overdraft or short loan to cover a gap.

\item hold less stock to free up cash.

\end{itemize}

Good \textbf{working capital management} keeps enough cash to stay safe without holding so much that money sits idle.

\section*{Costs and how we classify them}

A \textbf{cost} 成本 is money the business spends to make its product. We sort costs in several ways:

\begin{itemize}
\item \textbf{fixed costs} 固定成本 — costs that do not change with output, such as rent.

\item \textbf{variable costs} 可变成本 — costs that rise and fall with output, such as raw materials.

\item \textbf{direct costs} 直接成本 — costs clearly linked to one product, such as its materials.

\item \textbf{indirect costs} 间接成本 (overheads) — costs not linked to one product, such as manager salaries.

\item \textbf{marginal cost} 边际成本 — the cost of making one more unit.

\item \textbf{average cost} 平均成本 — the cost per unit, found by total cost ÷ number of units.

\end{itemize}

Cost information helps managers set prices, choose what to make, and control spending.

\section*{Break-even analysis}

\textbf{Contribution} 贡献 is how much each sale adds towards paying the fixed costs:

\[
\text{contribution per unit} = \text{selling price} - \text{variable cost per unit}
\]

\textbf{Break-even} 盈亏平衡 is the point where total revenue equals total costs, so profit is zero. The \textbf{break-even point} 盈亏平衡点 in units is:

\[
\text{break-even point} = \frac{\text{fixed costs}}{\text{contribution per unit}}
\]

The \textbf{margin of safety} 安全边际 is how far current sales are above the break-even point:

\[
\text{margin of safety} = \text{actual output} - \text{break-even output}
\]

A large margin of safety means sales can fall a lot before the firm makes a loss. Break-even analysis is quick and useful, but it has \textbf{limitations} 局限性: it assumes the selling price and costs stay the same, and that everything made is sold.

\section*{Budgets and variances}

A \textbf{budget} 预算 is a financial plan for the future — a target for income or spending. Budgets help a firm plan, control money, and check performance. Two types are:

\begin{itemize}
\item \textbf{incremental budget} 增量预算 — last year's budget changed by a small amount.

\item \textbf{zero-based budget} 零基预算 — every cost must be explained from zero each year.

\end{itemize}

A \textbf{variance} 差异 is the difference between the budgeted figure and the actual figure:

\[
\text{variance} = \text{actual figure} - \text{budgeted figure}
\]

\begin{itemize}
\item a \textbf{favourable variance} 有利差异 is better for profit (costs lower, or revenue higher, than planned).

\item an \textbf{adverse variance} 不利差异 is worse for profit (costs higher, or revenue lower, than planned).

\end{itemize}

Studying variances shows managers where the plan went wrong, so they can act.


\end{document}
